\section{Scalability, Performance and Resource Management}
\noindent\fbox{\parbox{\dimexpr\textwidth-2\fboxsep-2\fboxrule}{\small
\textbf{Tested vs.\ projected.} All numbers in this chapter labelled \emph{measured} were collected by us
on the development machine described in Section~8.6, with the real college timetable loaded. Numbers
labelled \emph{projected} are estimates and have not been tested on production hardware.}}

\subsection{Scalability Requirements}
\begin{table}[H]
\centering
\caption{Expected scale of one institute (projected)}
\label{tab:scale}
\small
\begin{tabularx}{\textwidth}{@{}Xp{6cm}@{}}
\toprule
\textbf{Quantity} & \textbf{Expected value}\\
\midrule
Registered users (students + staff) & 2{,}000--3{,}000\\
Batches / lab groups & 50--100 (real PDF: 46)\\
Class entries per official timetable & 500--1{,}000 (real PDF: 453 classes, 570 rows)\\
Timetable versions per semester & 3--10\\
Peak concurrent users (e.g.\ just after a timetable change) & 50--200\\
Peak request rate & 20--60 requests/s (a few requests per active user per minute)\\
Uploads & a few per week; each $\leq$ 15 MB, $\leq$ 30 pages\\
\bottomrule
\end{tabularx}
\end{table}

\subsection{Workload and Load Model}
The load test simulates students using the app at the same time:
\begin{itemize}
\item $N \in \{1, 2, 4, 8, 16, 32\}$ concurrent clients (threads), each logged in as a
      \emph{different} real student account (37 accounts, one per lab group).
\item Each client sends 30 search requests back-to-back with no think time (closed-loop, worst case),
      cycling through five realistic questions: ``What is my next class?'', ``What are my classes on
      Monday?'', ``Do I have classes after lunch on Tuesday?'', ``Which rooms are free at 10~am
      today?'' and ``Find all DBMS classes''.
\item Total requests per run $= 30N$ (30 to 960).
\item Every response is checked: HTTP 200 with \texttt{OK}/\texttt{NO\_MATCH}, or the expected 422
      clarification. (The M.Tech.\ page has no lunch row, so ``after lunch'' correctly asks for a time.)
\end{itemize}

\subsection{Resource Utilisation (measured)}
\begin{table}[H]
\centering
\caption{Measured resource utilisation}
\label{tab:resources}
\small
\begin{tabularx}{\textwidth}{@{}p{5.4cm}p{3cm}X@{}}
\toprule
\textbf{Resource} & \textbf{Value} & \textbf{How measured}\\
\midrule
Memory per API worker process (after load test) & 113--116 MB & \texttt{ps -o rss}, 4 workers\\
Peak memory, extraction of the 10-page PDF & 82 MB & \texttt{getrusage} max RSS\\
Database size (one timetable, 91 users, history) & 12 MB & \texttt{pg\_database\_size}\\
Entry table incl.\ indexes (570 rows) & 592 kB & \texttt{pg\_total\_relation\_size}\\
Original PDF stored & 420 kB & file size\\
CPU during 32-client test & $\approx$76\% of 4 cores & process CPU time / wall time\\
Web app bundle (JavaScript) & 3.5 MB uncompressed & exported bundle size\\
Network per search & $\approx$1--6 kB JSON & response size (depends on result count)\\
\bottomrule
\end{tabularx}
\end{table}
The KPI targets were under 300~MB per API worker and under 50~MB of database; both are met with a
large margin.

\subsection{Concurrent Processing}
Requests are handled by several independent Uvicorn worker \emph{processes} (4 in the test), each
with its own connection pool. Because the API is stateless, a request can go to any process. Within a
process, FastAPI runs synchronous endpoints in a thread pool, so slow database calls do not block other
requests. Extraction never runs in an API request: it is a job processed by the worker, so a large
OCR job cannot slow down searches. Shared state lives only in PostgreSQL and is protected by row locks
(Section~7.7).

\subsection{Scheduling and Synchronisation Analysis}
\begin{itemize}
\item \textbf{Request scheduling} is done by the operating system (processes) and Uvicorn's event
      loop. Under overload, requests queue at the socket. The rising p50 latency beyond 8 clients
      (Table~\ref{tab:loadresults}) is this queueing delay, since throughput stays flat at the CPU limit.
\item \textbf{Job scheduling} is FIFO by \texttt{created\_at}. Jobs are short (about 1.5~s for the
      real PDF), so FIFO gives low waiting times and needs no pre-emption. Round-Robin would only help
      with long jobs that need to share a worker, which our workload does not have.
\item \textbf{Synchronisation measurements:} the concurrency tests start 2 activations and 2
      promotions at the same moment. In every run exactly one wins and the other sees the committed
      result, with no deadlock and no lost update. The load test reported 0 inconsistent answers.
\end{itemize}

\subsection{Performance Measurement}
\begin{itemize}
\item \textbf{Environment:} Linux VM, 4 vCPU, 15~GB RAM; PostgreSQL 16 on the same machine; API with
      4 Uvicorn workers; the load generator also on the same machine (which takes CPU from the server,
      so results are conservative).
\item \textbf{Data:} the college's real timetable (453 classes), 46 batches, 52 faculty, 42 rooms,
      91 user accounts.
\item \textbf{Tools:} a Python script with \texttt{httpx} and threads; \texttt{time.perf\_counter}
      timestamps around each request; percentiles computed from all samples; \texttt{ps} for CPU and
      memory. For the load test only, the per-user rate limits were raised, because otherwise the
      limiter (60 searches/min per user) correctly rejects the synthetic burst.
\item \textbf{Single-user latency:} 50 sequential requests per question type after one warm-up request.
\end{itemize}

\subsection{Performance Results (measured)}
\begin{table}[H]
\centering
\caption{Single-user search latency per question type (50 requests each)}
\label{tab:latency}
\small
\setlength{\tabcolsep}{4pt}\begin{tabular}{@{}llrrr@{}}
\toprule
\textbf{Intent} & \textbf{Question} & \textbf{Results} & \textbf{p50 (ms)} & \textbf{p95 (ms)}\\
\midrule
Q10 Next class & What is my next class? & 1 & 42 & 52\\
Q02 Day timetable & What are my classes on Monday? & 5 & 44 & 52\\
Q06 Course search & Find all DBMS classes & 7 & 45 & 51\\
Q11 After lunch & Do I have classes after lunch on Tuesday? & 3 & 47 & 54\\
Q15 Professor schedule & What is my schedule on Monday? & 3 & 37 & 46\\
Q22 Free rooms & Which rooms are free at 10~am today? & 20 & 43 & 51\\
Q20 Room schedule & Show the schedule of room 508 on Monday & 8 & 41 & 49\\
\bottomrule
\end{tabular}
\end{table}

\begin{table}[H]
\centering
\caption{Load test with increasing concurrent clients (4 API workers)}
\label{tab:loadresults}
\small
\begin{tabular}{@{}rrrrrrr@{}}
\toprule
\textbf{Clients} & \textbf{Requests} & \textbf{Wrong} & \textbf{req/s} & \textbf{p50 (ms)} & \textbf{p95 (ms)} & \textbf{CPU (\%)}\\
\midrule
1 & 30 & 0 & 18.6 & 49 & 57 & 31\\
2 & 60 & 0 & 25.4 & 75 & 90 & 32\\
4 & 120 & 0 & 46.5 & 64 & 90 & 48\\
8 & 240 & 0 & 60.4 & 90 & 141 & 63\\
16 & 480 & 0 & 54.1 & 183 & 354 & 62\\
32 & 960 & 0 & 57.4 & 410 & 635 & 76\\
\bottomrule
\end{tabular}
\end{table}

\begin{figure}[H]
\centering
\begin{tikzpicture}
\begin{axis}[width=0.48\textwidth,height=6cm,xmode=log,log basis x=2,xtick={1,2,4,8,16,32},
  xticklabels={1,2,4,8,16,32},xlabel={Concurrent clients},ylabel={Throughput (req/s)},
  ymin=0,ymax=70,grid=major,title={\small Throughput}]
\addplot[thick,mark=*,color=spitblue] coordinates {(1,18.6)(2,25.4)(4,46.5)(8,60.4)(16,54.1)(32,57.4)};
\end{axis}
\end{tikzpicture}\hfill
\begin{tikzpicture}
\begin{axis}[width=0.48\textwidth,height=6cm,xmode=log,log basis x=2,xtick={1,2,4,8,16,32},
  xticklabels={1,2,4,8,16,32},xlabel={Concurrent clients},ylabel={Latency (ms)},
  ymin=0,grid=major,legend pos=north west,legend style={font=\scriptsize},title={\small Latency}]
\addplot[thick,mark=*,color=spitblue] coordinates {(1,49)(2,75)(4,64)(8,90)(16,183)(32,410)};
\addplot[thick,mark=square*,color=orange] coordinates {(1,57)(2,90)(4,90)(8,141)(16,354)(32,635)};
\legend{p50,p95}
\end{axis}
\end{tikzpicture}
\caption{Throughput and latency under increasing concurrent load (measured)}
\label{fig:load}
\end{figure}

\begin{table}[H]
\centering
\caption{Extraction performance on the real 10-page PDF (median of 5 runs)}
\label{tab:extractperf}
\small
\begin{tabular}{@{}lr@{}}
\toprule
\textbf{Metric} & \textbf{Value}\\
\midrule
Pages / rows produced & 10 / 570 (453 classes)\\
Extraction time, whole pipeline & 1.65 s (runs: 1.72, 1.65, 1.64, 1.69, 1.64 s)\\
Per page & $\approx$0.17 s\\
Peak memory & 82 MB\\
\bottomrule
\end{tabular}
\end{table}

\subsection{Bottleneck Identification}
\begin{enumerate}
\item \textbf{CPU of the API processes.} Throughput stops growing at about 55--60 req/s from 8 clients
      while CPU reaches 60--76\%, and the rest of the CPU is taken by PostgreSQL and the load generator
      on the same machine. Beyond that point extra clients only add queueing delay, so p50 rises
      roughly linearly with $N$.
\item \textbf{Per-request work:} each search resolves the context, loads the vocabulary of the
      timetable (about 200 labels) and runs one or two templates. Vocabulary loading is the largest
      repeated cost.
\item \textbf{Database round trips:} before optimisation a search issued about 50 SQL statements
      because related rows were loaded lazily (N+1 problem).
\end{enumerate}

\subsection{Optimisation and Improvement}
\begin{table}[H]
\centering
\caption{Optimisations performed (before/after)}
\label{tab:optim}
\small
\begin{tabularx}{\textwidth}{@{}p{4cm}p{3.2cm}p{3.2cm}X@{}}
\toprule
\textbf{Change} & \textbf{Before} & \textbf{After} & \textbf{Effect}\\
\midrule
Eager loading of course/staff/room/batch (N+1 fix) & $\approx$50 SQL statements / search & $\approx$25 & Halved round trips; latency now 37--47~ms\\
Batched vocabulary loading & 1 query per label & 1 query per entity type & Constant number of queries per search\\
Composite indexes led by \texttt{timetable\_id} & -- & index range scans & ``Now/next'' independent of other timetables' size\\
Intended-overlap rules in conflict detection & 542 reported overlaps & 2 & Removes false alarms (quality, not speed)\\
Extraction moved to background worker & -- & API never blocked by OCR & Search latency unaffected by uploads\\
\bottomrule
\end{tabularx}
\end{table}
Further options: cache the vocabulary per (timetable, \texttt{data\_version}); add more API processes
or machines (horizontal scaling); use a shared rate-limit store (Redis/PostgreSQL) for several machines.

\subsection{Cloud Deployment and Resource Utilisation}
\begin{itemize}
\item \textbf{Architecture:} Figure~\ref{fig:deploy}. One Render \emph{free} web service runs a Docker
      image (multi-stage build: Node builds the web app, then a Python image with Tesseract runs
      FastAPI, the worker thread and the static web app), plus a free managed PostgreSQL database.
\item \textbf{Configuration:} \texttt{render.yaml} (blueprint); the secret key is generated by Render,
      \texttt{DEMO\_PASSWORD} is entered at deploy time, production mode refuses the published
      development password, and demo data is seeded only into an empty database.
\item \textbf{Resources:} free instance with 512~MB RAM and a shared fraction of a CPU. Measured
      memory of one API process ($\approx$115~MB) plus extraction peak ($\approx$82~MB) fits within it.
\item \textbf{Limits:} the service sleeps after 15 minutes idle (first request then takes about a
      minute); the free database expires after 30 days; 750 free instance hours per month~\cite{render}.
\item \textbf{Scaling approach (projected):} a paid always-on instance with 2+ API processes and a
      separate worker; managed PostgreSQL with backups; horizontal scaling behind the platform's load
      balancer, because the API is stateless.
\end{itemize}

\subsection{Scalability Limitations}
\begin{itemize}
\item \textbf{Maximum tested workload:} 32 concurrent closed-loop clients (960 requests) with 0 wrong
      answers and 57 req/s, on a 4-core development VM shared with the database and the load generator.
\item \textbf{Projection (not tested):} with think time of about 10~s between requests, 57 req/s
      corresponds to several hundred simultaneously active users on similar hardware. This is an
      estimate only; a real capacity test must run on the deployment hardware.
\item The free cloud tier is much smaller (shared CPU); it is suitable for a demo, not for the whole college.
\item The rate limiter is per process. With several machines it must move to a shared store.
\item OCR of large scanned PDFs is CPU-heavy; with many simultaneous uploads more workers are needed.
\end{itemize}

\newpage
